\documentclass[11pt]{article}
\usepackage{amsmath, amssymb, amsthm}
\usepackage[shortlabels]{enumitem}
\usepackage[margin=1in]{geometry}

\newtheorem{theorem}{Theorem}
\newtheorem{definition}[theorem]{Definition}
\theoremstyle{definition}
\newtheorem{exercise}{Exercise}

\title{Worksheet 4: Matrix Operations and Row Reduction}
\author{Math 22a --- Linear Algebra}
\date{Fall 2025}

\begin{document}
\maketitle

\begin{definition}
An $m \times n$ \textbf{matrix} is a rectangular array of numbers with $m$ rows and $n$ columns:
\[
  A = \begin{pmatrix}
    a_{11} & a_{12} & \cdots & a_{1n} \\
    a_{21} & a_{22} & \cdots & a_{2n} \\
    \vdots & \vdots & \ddots & \vdots \\
    a_{m1} & a_{m2} & \cdots & a_{mn}
  \end{pmatrix}.
\]
We write $A = (a_{ij})$ where $a_{ij}$ denotes the entry in row $i$, column $j$.
\end{definition}

\begin{exercise}
Perform the indicated matrix operations, or explain why the operation is not defined.
\begin{enumerate}[(a)]
  \item $\begin{pmatrix} 2 & -1 \\ 0 & 3 \end{pmatrix}
         + \begin{pmatrix} 1 & 4 \\ -2 & 5 \end{pmatrix}$

  \item $3 \begin{pmatrix} 1 & 0 & -2 \\ 4 & 1 & 3 \end{pmatrix}
         - \begin{pmatrix} 2 & -1 & 0 \\ 0 & 3 & -6 \end{pmatrix}$

  \item $\begin{pmatrix} 1 & 2 \\ 3 & 4 \end{pmatrix}
         \begin{pmatrix} 0 & 1 \\ 1 & 0 \end{pmatrix}$

  \item $\begin{pmatrix} 1 & 2 & 3 \end{pmatrix}
         \begin{pmatrix} 4 \\ 5 \\ 6 \end{pmatrix}$
\end{enumerate}
\end{exercise}

\begin{exercise}
Use Gaussian elimination to reduce each augmented matrix to row echelon form, then solve the corresponding system.
\begin{enumerate}[(a)]
  \item $\left(\begin{array}{ccc|c}
    1 & 2 & -1 & 3 \\
    2 & 5 & 0  & 7 \\
    -1 & 1 & 3 & 0
  \end{array}\right)$

  \item $\left(\begin{array}{ccc|c}
    1 & -1 & 2  & 4 \\
    3 & -3 & 6  & 12 \\
    2 & 1  & -1 & 2
  \end{array}\right)$
\end{enumerate}
\end{exercise}

\begin{exercise}
For each matrix below, find the reduced row echelon form and determine the rank.
\begin{enumerate}[(a)]
  \item $A = \begin{pmatrix} 1 & 3 & 5 & 7 \\ 2 & 4 & 6 & 8 \\ 0 & 1 & 1 & 1 \end{pmatrix}$

  \item $B = \begin{pmatrix} 1 & 0 & 2 \\ 0 & 1 & -1 \\ 2 & 1 & 3 \\ 1 & -1 & 3 \end{pmatrix}$
\end{enumerate}
\end{exercise}

\begin{exercise}
Determine whether each set of vectors is linearly independent.
\begin{enumerate}[(a)]
  \item $\mathbf{v}_1 = \begin{pmatrix} 1 \\ 2 \\ 3 \end{pmatrix}, \quad
         \mathbf{v}_2 = \begin{pmatrix} 4 \\ 5 \\ 6 \end{pmatrix}, \quad
         \mathbf{v}_3 = \begin{pmatrix} 7 \\ 8 \\ 9 \end{pmatrix}$

  \item $\mathbf{u}_1 = \begin{pmatrix} 1 \\ 0 \\ 1 \end{pmatrix}, \quad
         \mathbf{u}_2 = \begin{pmatrix} 0 \\ 1 \\ 1 \end{pmatrix}, \quad
         \mathbf{u}_3 = \begin{pmatrix} 1 \\ 1 \\ 0 \end{pmatrix}$
\end{enumerate}
\end{exercise}

\begin{exercise}
Let $T \colon \mathbb{R}^2 \to \mathbb{R}^2$ be the linear transformation defined by
\[
  T\begin{pmatrix} x \\ y \end{pmatrix} = \begin{pmatrix} 2x - y \\ x + 3y \end{pmatrix}.
\]
\begin{enumerate}[(a)]
  \item Write the standard matrix $A$ for $T$.
  \item Compute $T\begin{pmatrix} 1 \\ -1 \end{pmatrix}$.
  \item Is $T$ one-to-one? Is $T$ onto? Justify your answers.
  \item Find $\det(A)$ and use it to determine whether $A$ is invertible.
\end{enumerate}
\end{exercise}

\begin{theorem}[Invertible Matrix Theorem --- Partial List]
Let $A$ be an $n \times n$ matrix. The following are equivalent:
\begin{enumerate}[(i)]
  \item $A$ is invertible.
  \item $A$ is row equivalent to $I_n$.
  \item $A$ has $n$ pivot positions.
  \item The equation $A\mathbf{x} = \mathbf{0}$ has only the trivial solution.
  \item The columns of $A$ are linearly independent.
  \item $\det(A) \neq 0$.
\end{enumerate}
\end{theorem}

\end{document}
